O módulo de resistência à flexão é dado por:

\begin{equation}
  W_f = \frac{M_J}{\sigma_{\text{adm}}}
  \label{eq:w_f_perna}
\end{equation}

Onde:

\begin{itemize}
  \item $W_f$: módulo de resistência à flexão em \si{\meter^3}.
  \item $M_J$: momento de reação no ponto J em \si{\newton\meter} (Tabela \ref{tab:resultados_vp1}).
  \item $\sigma_{\text{adm}}$: tensão de flexão admissível em \si{\pascal} \eqref{eq:tensao_admissivel_aco_a36}.
\end{itemize}

\vspace{1em}

\begin{align*}
  W_f & = \frac{\num{48398.89}}{\num{187.97e6}} \\
  W_f & = \boxed{\SI{2.575e-4}{\meter^3}}
\end{align*}

\begin{align*}
  W_{zz} & = W_f \cdot 10^9                           \\
  W_{zz} & = \num{2.575e-4} \cdot 10^9                \\
  W_{zz} & = \boxed{\SI{2.575e5}{\milli\meter\cubed}}
\end{align*}

Com base no Projeto Preliminar (Figura \ref{fig:alturas_vigas}), foi adotada a seção de $h = \SI{270}{\milli\meter}$.

\begin{table}[H]
  \centering
  \caption{Dimensões da seção da perna do pórtico}
  \label{tab:dimensoes_perna_portico}
  \begin{tabularx}{\textwidth}{>{\centering\arraybackslash}X>{\centering\arraybackslash}X>{\centering\arraybackslash}X}
    \toprule
    Parâmetro                      & Símbolo  & Valor                              \\
    \midrule
    Altura                         & $h$      & \SI{270}{\milli\meter}             \\
    Largura                        & $b$      & \SI{90}{\milli\meter}              \\
    Espessura                      & $e$      & \SI{6.3}{\milli\meter}             \\
    Módulo de resistência à flexão & $W_{zz}$ & \SI{280e3}{\milli\meter\cubed}     \\
    Massa por metro quadrado       & $m$      & \SI{49.89}{\kilo\gram\per\meter^2} \\
    Largura interna                & $b_1$    & \SI{64.8}{\milli\meter}            \\
    Altura interna                 & $h_1$    & \SI{257.4}{\milli\meter}           \\
    \bottomrule
  \end{tabularx}
\end{table}
